\begin{abstract}
Thermal cameras make nighttime pedestrian segmentation reliable, but their cost and bandwidth keep RGB-Thermal (RGB-T) perception out of low-cost robots and vehicles. We ask a budget question: how many thermal pixels are actually needed, and where should they be used? On MFNet, we first show that the thermal teacher's nighttime pedestrian advantage cannot be transported to an RGB-only student by logit distillation, modality dropout, or thermal hallucination (night-person TCE $\approx -0.08$): the evidence is modality-borne. We then sweep the thermal sensor from $60\times80$ to $240\times320$ and find a clear knee at $120\times160$: $1/16$ of the RGB pixel count recovers 81\% of the full-resolution teacher's nighttime pedestrian gain. Used globally, however, low-resolution thermal loses on eight other classes. We therefore introduce a person-aware reliability gate: a small U-Net, trained on a held-out split of the training set with hard-condition augmentation, decides per pixel when to defer from an RGB expert to the $120\times160$ thermal expert, and a structural constraint restricts thermal overrides to pedestrians. With three seeds and a frozen 392-image test, the gate raises overall mIoU from 56.38 to $57.57\pm0.08$ and nighttime pedestrian IoU from 45.24 to $61.36\pm0.39$; all other classes change by at most 0.37 IoU. Under synthetic low-light, blur, and noise, RGB collapses to 47.95 mIoU while the gate holds 51.10. Thermal sensing should be budgeted and class-routed, not distilled or uniformly deployed.
\end{abstract}
